\scanpage{139}
\deleted{\unsure{du fait d'homogénéité|à part d'homogénéité} \unsure{il vient,|évidente,} qui permet}

Remplaçons dans cette formule $A$ par $\lambda A$, $B$ par $\mu B$ ($\lambda$ et $\mu$ nombres $>0$) on obtient.

\[
\deleted{\unsure{\lambda^p\mu^q=\frac{\lambda^p\mu^q}{r}S_r(AB)}}
\inserted{\frac{(\lambda\mu)^r}{r}S_r(AB)}
\leq\frac{\lambda^p}{p}S_p(A)+\frac{\mu^q}{q}S_q(B)
\]

\deleted{Exprimons que cette inégalité doit être \unsure{vérifiée} quels que soient $x,y\unsure{\geq|>}0$.}

\margin{\unsure{Ici lemme que} la condition du \unsure{n°} implique $c^{\frac1r}\leq a^{\frac1p}b^{\frac1q}$ i.e \uuline{$c\leq a^\alpha b^\beta$}}

Pour simplifier, posons $a=S_p(A)$, $b=S_q(B)$, $c=S_{\unsure{p|r}}(C)$, on \unsure{aura}

\[
c\leq\mathop{\mathrm{Inf}}\limits_{\substack{\lambda\geq0\\\mu\geq0}}
\frac{a\frac{\lambda^p}{p}+b\frac{\mu^q}{q}}
{\substack{\deleted{\unsure{(\lambda\mu)^r}}\\\deleted{\unsure{\frac{\lambda^p\mu^q}{r}}}}\frac{(\lambda\mu)^r}{r}}
\]

\deleted{Posons $\unsure{x^p=\lambda,y^q}$}\deleted{\inserted{$x=\lambda^p,\ y=\mu^q$,}}

Posons $x=\lambda^p$, $y=\mu^q$, $\alpha=\frac rp$, $\beta=\frac rq$ \unsure{.|,} d'où \unsure{,} \uline{$\frac rp+\frac rq=1$},

On a
\[
c\leq\mathop{\mathrm{Inf}}\limits_{\substack{x>0\\y>0}}
\frac{\alpha a\deleted{\unsure{\frac rp}}x+\deleted{\unsure{\frac rq}}\beta by}{x^\alpha y^\beta}
=\mathop{\mathrm{Inf}}\limits_{\substack{x\geq0\\y\geq0\\x^\alpha y^\beta=1}}\alpha a x+\beta by
\]

Or on vérifie aussitôt que \unsure{la|le} \deleted{\unsure{minimum}}\inserted{borne \unsure{inférieure}} du second membre est atteinte (car l'en\unsure{s.} \unsure{des} points $(x\unsure{,|;}y)$ tels que $x\geq0,y\geq0$ $ax+by\leq M$ -- valeur donnée -- est compact) elle est donc atteinte en un point $(x,y)$ tel que les équations \deleted{$\unsure{a\,dx+b\,dy=0}$}\inserted{$d(\alpha ax\unsure{-|+}\beta by)=0$}, $d(x^\alpha y^\beta)=0$ \deleted{soient compatibles} \inserted{aient une solution non nulle en} $dx,dy$ ;

\unsure{or} $d(x^\alpha y^\beta)=\alpha x^{\alpha-1}y^\beta\,dx+\beta y^{\beta-1}x^\alpha\,dy$,\margin{\unsure{2|?}}
\deleted{i.e\unsure{.} tel que $\unsure{\det(\alpha a\quad\beta b\quad\alpha a}$}
la condition d'extremum est donc $\frac{\alpha a}{\alpha x^{\alpha-1}y}=\frac{\beta b}{\beta x^\alpha y^{\beta-1}}$

soit $ax=by$ \deleted{\unsure{$=$}}, \deleted{\unsure{et comme on a}} \inserted{d'où $ax=by=\frac{\alpha ax+\beta by}{\alpha+\beta}$} ; \deleted{$\unsure{\alpha ax+by=\frac{\alpha ax}{\unsure{\alpha|2}}\unsure{=|+}\frac{\alpha ax+\beta by}{\unsure{\alpha|2}}}$} Comme $\alpha+\beta=1$, le minimum cherché sera alors \deleted{\unsure{$ax=by$|$ax+by$|$\alpha x+b$}} égal à $ax=by$ ; compte tenu de la \unsure{condi}

